% Generated draft: requires mathematical and source review.
\documentclass[11pt]{article}
\usepackage[margin=25mm]{geometry}
\usepackage{fontspec}
\setmainfont{Latin Modern Roman}
\usepackage{amsmath,amssymb,mathtools,unicode-math}
\setmathfont{Latin Modern Math}
\usepackage{enumitem,needspace,xurl,hyperref}
\setlength{\parindent}{0pt}
\setlength{\parskip}{5pt}
\setlength{\emergencystretch}{3em}
\newcommand{\N}{\mathbb{N}}
\newcommand{\Z}{\mathbb{Z}}
\newcommand{\Q}{\mathbb{Q}}
\newcommand{\R}{\mathbb{R}}
\newcommand{\C}{\mathbb{C}}
\newcommand{\F}{\mathbb{F}}
\newcommand{\A}{\mathbb{A}}
\newcommand{\PP}{\mathbb{P}}
\newcommand{\E}{\mathbb{E}}
\newcommand{\eps}{\varepsilon}
\newcommand{\dd}{\,\mathrm{d}}
\begin{document}
\section*{UnsolvedMath ID: 700557}
\textbf{Source problem code:} catalogue entry
\subsection{Definitions and mathematical statement}
A group \(G\) is simple if its only normal subgroups are \(\{1\}\) and \(G\). For \(c\in G\), its conjugacy class is
\[
C_G(c)=\{xcx^{-1}:x\in G\}.
\]
Here \(C_G(c)\) denotes a conjugacy class, not a centralizer. For a subset \(C\subseteq G\), define its setwise square by
\[
CC=\{ab:a,b\in C\}.
\]
\textbf{Conjecture (J. G. Thompson; proposed in the cited source by V. D. Mazurov).} For every finite simple non-abelian group \(G\), there exists \(c\in G\) such that
\[
G=C_G(c)C_G(c).
\]
Equivalently, for every such \(G\), there exists \(c\in G\) such that for every \(g\in G\) there exist \(x,y\in G\) with
\[
g=(xcx^{-1})(ycy^{-1}).
\]
The supplied registry labels this problem open; that label is supplied source metadata and has not been independently verified.
\subsection{Short English statement}
Every finite simple non-abelian group is conjectured to have a single conjugacy class such that every group element is a product of two elements of that class. The statement is attributed to J. G. Thompson, with V. D. Mazurov recorded as the proposer.
\subsection{Sources}
\begin{itemize}
\item Supplied local UnsolvedMath registry record, permanent ID \(700557\), exact source code \texttt{catalogue entry}, titled \emph{Kourovka Notebook Problem 9.24}. The supplied status is \texttt{open}.
\item \emph{The Kourovka Notebook (Unsolved Problems in Group Theory)}, original problem 9.24, 9th issue (1984); also identified in the supplied record with the 21st issue edition (2026), arXiv:1401.0300v48, \url{https://arxiv.org/abs/1401.0300v48}. The supplied collection description names E. I. Khukhro and V. D. Mazurov as editors. This reference is reproduced from the supplied record and has not been independently consulted.
\item Attribution: the statement names J. G. Thompson; the proposer and background fields name V. D. Mazurov and give the proposal year as 1984.
\end{itemize}
\end{document}
